\pagenumbering{arabic}
\chapter{INTRODUCTION}

\section{Background}
Machine learning models have become deeply integrated into everyday applications, from personalized content recommendations to critical tasks like medical diagnosis and financial fraud detection. These models learn patterns from vast amounts of user data, improving their accuracy and usefulness over time. However, this dependence on user data creates a fundamental tension with privacy rights. Regulations such as the General Data Protection Regulation (GDPR) in Europe and the California Consumer Privacy Act (CCPA) in the United States now establish clear legal requirements: users have the right to request deletion of their personal data from systems that process it. The technical challenge lies in how to honor such deletion requests when the data has already been incorporated into a trained machine learning model. Traditional approaches require complete model retraining on the remaining dataset, a process that becomes prohibitively expensive as model complexity and dataset sizes grow. This gap between legal requirements and technical capability has given rise to the field of machine unlearning, which aims to efficiently remove the influence of specific data points from trained models without full retraining.

\section{Gap Identification}
Despite growing research interest in machine unlearning, the practical adoption of these techniques remains limited. Academic papers propose sophisticated algorithms, but they often remain as theoretical prototypes or are tightly coupled to specific model architectures. Open-source libraries that do exist tend to focus on narrow use cases or require significant expertise to integrate into existing machine learning pipelines. For most developers and organizations, there is no straightforward, production-ready solution that allows them to build privacy-compliant systems with unlearning capabilities from day one. This project addresses that gap by developing PyUnlearn, a practical SDK and API platform that abstracts away the complexity of machine unlearning, making it accessible to developers working on real-world applications that must comply with privacy regulations.

\section{Motivation}
The motivation for this project stems from a simple but increasingly urgent question: as users entrust their data to machine learning systems, can they also trust that their presence can be removed when requested? Current practices often treat data Qdeletion as a database operation, ignoring the fact that the model itself retains implicit knowledge of the deleted data. This disconnect between what users expect and what systems actually deliver creates legal risk for organizations and erodes user trust. Rather than treating privacy as an afterthought or a compliance checkbox, we believe that unlearning capabilities should be built directly into the machine learning development workflow. PyUnlearn represents an effort to make this possible, providing developers with tools that integrate naturally with existing workflows while ensuring that privacy guarantees are maintained throughout the model lifecycle.

\section{Objectives}
The project aims to achieve the following objectives:
\begin{itemize}
\item Develop an efficient selective unlearning framework that allows developers to integrate unlearning capabilities into their machine learning pipelines, with the flexibility to run locally or leverage cloud infrastructure for training and unlearning operations.
\item Ensure that unlearning operations provide meaningful privacy guarantees while preserving model performance on retained data, balancing the trade-off between computational efficiency and model accuracy.
\item Deliver a practical solution that enables compliance with GDPR, CCPA, and similar privacy regulations by providing auditable unlearning operations, persistent metadata tracking, and idempotent deletion semantics suitable for production deployment.
\end{itemize}

\section{Scope and Limitations}
This project focuses on implementing a practical machine unlearning SDK based on the SISA (Sharded, Isolated, Sliced, Aggregated) framework, targeting tabular datasets and machine learning models compatible with the scikit-learn ecosystem. The system supports both local deployment and cloud-enabled monitoring through an optional dashboard. While the approach demonstrates significant efficiency gains over full retraining, current limitations include the lack of support for deep learning architectures such as convolutional neural networks, reliance on a fixed sharding strategy that may not adapt optimally to all dataset distributions, and the absence of formal certification guarantees for the unlearning process. These limitations are acknowledged as areas for future extension beyond the scope of the current project.